% Fragmento conservado; véase MANIFIESTO_ANTECEDENTES_LECTURA.json.
\subsection{Symétrie spéculaire des régions et double réalisation}
\label{mono-ampl:regiones}

L'échange spéculaire agit sur le bloc régional
\(B=(u^\pi,u^e,u^\varphi)\in(\mathbb F_3^6)^3\) par
\[
 \mathfrak cB=(u^e,u^\pi,u^\varphi).
\]
Son axe fixe et son agrégat sont respectivement \(u^\varphi\) et
\(s=u^\pi+u^e\). Cette fixation concerne l'involution dans une fibre ;
les deux coordonnées peuvent évoluer pendant le transport nonadique.
Avec \(q=u^\pi+u^e+u^\varphi\),
\(a=u^\pi+u^e-u^\varphi\) et \(d=u^\pi-u^e\), on retrouve
\[
 u^\varphi=2(q-a),\qquad s=q-u^\varphi,\qquad
 u^\pi=2(s+d),\qquad u^e=2(s-d).
\]
Les égalités sont composante par composante dans \(\mathbb F_3\), où
\(2^{-1}=2\). Le miroir conserve \(q,a,s\) et inverse \(d\) ; cette
coordonnée antisymétrique individualise la répartition que l'agrégat conserve
sans la distinguer. L'inversion \(\mathcal I\) de l'indice temporel et
l'involution \(\mathfrak c\) des canaux agissent donc sur des
données différentes.

Le retour de phase possède un témoin explicite :
\[
 \begin{aligned}
 B_1&=(012222\mid121221\mid112202),\\
 B_{10}&=(101222\mid112221\mid112002).
 \end{aligned}
\]
Les deux positions ont la phase \(1\) ; le bloc, l'axe et la mémoire ont
changé. De même, la synchronisation des dénombrements aux phases \(8\) et
\(9\) annule leurs différences relatives, tandis que la composante diagonale
passe de \((59,59,59)\) à \((43,43,43)\). L'égalité d'un observable
relatif n'épuise pas l'état qui le détermine.

La double réalisation conserve cette architecture. Dans la section 8.1 de l'Article I (provenance du corridor exceptionnel, conservée dans le registre technique), le même préfixe détermine d'une part des cylindres emboîtés et d'autre part une suite de mots codés avec leurs repères. La première application passe à la limite par contraction des diamètres ; la seconde satisfait \(r_{n,m}Q_n=Q_mp_{n,m}\) et définit \(Q_\infty\). La figure \ref{mono-ampl:fig-regional} anticipe ces deux applications. Le parcours d'incidence ultérieur construit des codes, des réseaux et la réalisation de Moonshine dans leurs domaines propres ; le groupe d'automorphismes agit sur cette réalisation, et non sur les chiffres par une identification implicite.

% Fragmento conservado; véase MANIFIESTO_ANTECEDENTES_LECTURA.json.
% Figura acumulativa; destino figures/monodromia_regional_ampliacion.tex.
\begin{figure}[!htbp]
\centering
\begin{tikzpicture}[x=1cm,y=1cm,>=Latex,font=\small,
  flow/.style={->,line width=.55pt},
  ann/.style={align=center,font=\footnotesize}]
\node[font=\bfseries] at (6.9,9.2)
  {Connexion nonadique et coordonnées régionales};
\begin{scope}[shift={(3.4,6.8)}]
 \foreach \k in {1,...,9}{
  \pgfmathsetmacro{\ang}{90-40*(\k-1)}
  \coordinate (p\k) at (\ang:1.52);
  \fill (p\k) circle (1.4pt);
  \node[fill=white,inner sep=1.1pt,font=\footnotesize]
    at (\ang:1.83) {\(\k\)};
 }
 \foreach \k in {1,...,8}{
  \pgfmathtruncatemacro{\next}{\k+1}
  \draw[flow,shorten >=3pt,shorten <=3pt] (p\k)--(p\next);
 }
 \draw[flow,shorten >=3pt,shorten <=3pt] (p9)--(p1);
 \node[ann] at (0,.1) {\(g_k\in C_9\)\\\(\Gamma_{9,k}\)};
 \node[ann] at (0,-2.23) {Un tour : \(g\mapsto g\), \(m\mapsto m+1\)};
\end{scope}
\node[ann,text width=5.25cm] at (10.1,7.75)
 {\(B_k=(u_k^\pi,u_k^e,u_k^\varphi)\)\\
 bloc régional dans une fibre};
\node at (8.55,6.5) {\(u^\pi\)};
\node at (11.65,6.5) {\(u^e\)};
\draw[<->,line width=.65pt] (8.95,6.5)--node[above]{\(\mathfrak c\)}(11.25,6.5);
\draw[densely dashed,gray] (10.1,5.72)--(10.1,7.23);
\node[fill=white,inner sep=2pt] at (10.1,5.6){\(u^\varphi\)};
\node[ann] at (10.1,4.6)
 {Axe \(u^\varphi\) et agrégat \(u^\pi+u^e\)\\
 invariants sous \(\mathfrak c\)};
\draw[flow] (5.45,6.8)--(7.25,6.8);
\draw[gray!60,line width=.3pt] (.55,4.02)--(13.25,4.02);
\node[ann,text width=12cm] (hist) at (6.9,3.47)
 {Histoire compatible : préfixes, cylindres, orientation, quotients,
 repères d'incidence et mémoire};
\coordinate (fork) at (6.9,2.85);
\draw[flow] (hist.south)--(fork);
\draw[flow] (fork)--(3.6,2.2);
\draw[flow] (fork)--(10.3,2.2);
\node[ann,text width=5.6cm] at (3.6,1.65)
 {Réalisation archimédienne\\
 \(I_{n+1}\subseteq I_n,\quad |I_n|=729^{-n}\)\\
 limite des préfixes régionaux};
\node[ann,text width=5.6cm] at (10.3,1.65)
 {Réalisation incidencielle\\
 \(r_{n,m}Q_n=Q_mp_{n,m}\)\\
 mots codés et repères compatibles};
\node[ann,text width=5.6cm] at (3.6,.42)
 {\(\pi,\varphi,e\) ; détermination dodécaphasique de \(\alpha\)};
\node[ann,text width=5.6cm] at (10.3,.42)
 {Witt--Golay--Leech--\(V^\natural\)\\
 automorphismes de la réalisation exceptionnelle};
\draw[flow] (3.6,1.03)--(3.6,.78);
\draw[flow] (10.3,1.03)--(10.3,.78);
\end{tikzpicture}
\caption{\fontsize{9}{11}\selectfont Neuf phases d'une connexion et deux réalisations de l'histoire
compatible. Le cycle représente la phase ; chaque tour incrémente le compteur.
Le diagramme régional représente l'involution \(\mathfrak c\), qui
échange clôture et propagation et fixe l'axe d'auto-échelle dans une fibre.
Cette invariance n'impose pas que l'axe reste constant pendant le transport.
Les deux applications inférieures utilisent le même préfixe : l'une détermine
sa limite numérique et l'autre conserve les mots et les cadres d'incidence. Les
applications de cette seconde réalisation sont spécifiées dans la section
8.1 de l'Article I (provenance du corridor exceptionnel, conservée dans le registre technique) ; le parcours exceptionnel est développé ensuite.}
\label{mono-ampl:fig-regional}
\end{figure}


